\begin{equation}
\mathcal O\!\left(K(P_{\bm\theta}+MN+M^3)\right), \label{eq:onlinecomplex}
\end{equation}
plus one final forward evaluation of the same order. For a batch of $B$ samples,
an offline training iteration scales as
$\mathcal O(BK(P_{\bm\theta}+MN+M^3))$. Because the implemented outer update is
first-order, it does not store a full second-order graph through all $K$ inner
steps, reducing memory and derivative cost relative to full MAML. This is the
relevant complexity for \emph{prediction}: one compact network plus a fixed,
short RIS adaptation, rather than solving a fresh long alternating-optimization
problem from scratch for every channel.

%==============================================================================
\section{Numerical Results and Analysis}
Table~\ref{tab:sim} reports the parameters used by the fair NOMA-RIS simulator
and the final learning/evaluation scripts. The channel generator uses
$10{,}000$ realizations with a 90/10 train/validation split. The node geometry is
specified in the simulator by link distances rather than unique Cartesian
coordinates; reporting these distances therefore reproduces the propagation
model without inventing coordinates that are not present in the code.

\begin{table}[t]
\centering
\caption{System, channel, and learning parameters.}
\label{tab:sim}
\scriptsize
\begin{tabular}{@{}ll@{}}
\toprule
Parameter & Value \\ \midrule
BS antennas $M$ & 4 \\
RIS elements $N$ & 8 nominal; $\{8,16,32,64\}$ in RIS-size sweep \\
Power-sweep RIS size & $N=16$ \\
Monte-Carlo realizations & 10,000 (90\% train, 10\% validation) \\
BS--RIS distance $d_{BR}$ & 8 m \\
RIS--near distance $d_{RD_n}$ & 3 m \\
RIS--far distance $d_{RD_f}$ & 25 m \\
RIS--target / target--RIS & $d_{RT}=d_{TR}=6$ m \\
Direct BS--near / BS--far refs. & $d_{BD_n}=10$ m, $d_{BD_f}=25$ m \\
Reference path gain $PL_0$ & $10^{-3}$ ($-30$ dB) \\
Path-loss exponents & $2.2,2.8,2.8,2.3,2.3,3.8,3.8$ \\
Nakagami-$m$ parameters & $m_{BR}=2$, $m_{RD_n}=m_{RD_f}=1$; \\
 & $m_{RT}=m_{TR}=3$, $m_{BD_n}=m_{BD_f}=1$ \\
Target gain $\beta_T$ & 40 dB \\
RIS amplitude $\beta_i$ / efficiency & 1 / 1 (ideal passive RIS) \\
Bandwidth / temperature / NF & 1 MHz / 290 K / 6 dB \\
Generator power $P_{\rm tot}$ & 40 dBm; trainer permits power override \\
RIS-size figure operating power & 15 dBm \\
Power sweep & $5,10,15,20,25$ dBm \\
Initial NOMA split $(a_n,a_f,a_T)$ & $(0.20,0.55,0.25)$ \\
Communication beam weights & far/near $=0.7/0.3$ \\
Sensing leakage penalty $\lambda_{\rm null}$ & 0.5 \\
Objective weights $(\omega_c,\omega_s)$ & $(0.7,0.3)$ \\
NOMA cluster/sensing QoS & $R_{\rm th,c}=21.7$, $R_{\rm th,s}=2.9$ bit/s/Hz \\
Individual training floors & $R_{\rm th,n}=R_{\rm th,f}=1$ bit/s/Hz \\
Penalty weights $(\lambda_c,\lambda_s,\lambda_n,\lambda_f)$ & $(2,50,15,20)$ \\
Power network & 9--256--256--256--3, softmax output \\
NOMA-RIS result checkpoints & batch 64, lr $10^{-4}$, 40 epochs, $K=10$ \\
OMA-RIS result checkpoints & batch 64, lr $10^{-4}$, 40 epochs, $K=5$ \\
No-RIS result checkpoints & batch 128, lr $5\times10^{-4}$, 60 epochs \\
Inner phase optimizer & Adam, learning rate 0.1 \\
\bottomrule
\end{tabular}
\end{table}

The aggregate NOMA threshold in Table~\ref{tab:sim} comes from the fair NOMA
dataset, where it was chosen to keep the cluster QoS active; the individual
near/far floors are the defaults used by the training scripts. The OMA fair
dataset uses its own calibrated aggregate/sensing thresholds
($R_{\rm th,c}=12$ and $R_{\rm th,s}=2$ bit/s/Hz). The service-rate outage
thresholds stored by the MATLAB generators are diagnostic quantities and are
not the same as the optimization penalties in~\eqref{eq:sampleloss}.

The implementation records per-iteration loss and weighted rate for convergence
monitoring, while model selection uses a held-out validation score. The present
analysis therefore relies on the fixed-step inner optimizer, validation-based
checkpointing, and the convergence limitations discussed in
Section~III-E rather than asserting a closed-form global guarantee.

Fig.~\ref{fig:training} shows the MAML meta-loss~\eqref{eq:sampleloss} over training
iterations for several batch/learning-rate settings. The loss decreases smoothly
and the weighted sum-rate climbs steadily, with no epoch-one collapse,
indicating stable second-order meta-optimization; the steady-state panel
confirms continued, generalizing improvement.

\begin{figure}[t]
  \centering
  \includegraphics[width=\columnwidth]{fig_training.eps}
  \caption{MAML meta-loss versus training iteration over a batch/learning-rate
  sweep ($M{=}4$, $N{=}8$). (a)~Full trajectory; (b)~steady state
  ($\text{iter}\!\ge\!200$).}
  \label{fig:training}
\end{figure}

Fig.~\ref{fig:cdf} plots the empirical CDF of the weighted sum-rate $\Rdl$ over
the test set for the ablation. Joint optimization of phase and power
(\textsf{fullOptimized}) stochastically dominates the partial designs, raising
the median from $\sim\!5.0$ to $\sim\!6.7$~bps/Hz; phase optimization accounts
for most of the rate gain, while the learned power split adds a further,
QoS-critical improvement. This is reflected in Fig.~\ref{fig:qos}: the joint
QoS-violation rate falls monotonically from $29.4\%$ (no optimization) to
$0.1\%$ (proposed), with both the communication ($q_c$) and sensing ($q_s$)
constraints satisfied almost surely.

\begin{figure}[t]
  \centering
  \includegraphics[width=\columnwidth]{fig_cdf.eps}
  \caption{CDF of the weighted sum-rate $\Rdl$ over the test set for the five
  ablation configurations.}
  \label{fig:cdf}
\end{figure}

\begin{figure}[t]
  \centering
  \includegraphics[width=\columnwidth]{fig_qos.eps}
  \caption{QoS-violation fraction (lower is better). The joint violation rate
  drops from $29.4\%$ to $0.1\%$ along the ablation.}
  \label{fig:qos}
\end{figure}

% --- OPTIONAL extra ablation figures (uncomment to include) ------------------
% \begin{figure}[t]\centering
%   \includegraphics[width=\columnwidth]{fig_improvement.jpg}
%   \caption{Per-metric improvement over the unoptimized baseline.}\label{fig:improvement}\end{figure}
% \begin{figure}[t]\centering
%   \includegraphics[width=\columnwidth]{fig_perleg.jpg}
%   \caption{Per-leg rate CDFs ($R_n$, $R_f$, $R_s$).}\label{fig:perleg}\end{figure}

Fig.~\ref{fig:rateN} reports $\Rsum$ versus the RIS size
$N\in\{8,16,32,64\}$ at $P=15$~dBm. The rate grows monotonically with $N$ for
both access schemes—from $6.9$ to $11.4$~bps/Hz for NOMA and from $4.1$ to
$7.3$~bps/Hz for OMA—and NOMA retains an $\sim\!65\%$ advantage throughout
showing that the passive aperture and the access scheme contribute complementary
gains.

\begin{figure}[t]
  \centering
  \includegraphics[width=\columnwidth]{fig_rate_N.eps}
  \caption{Weighted sum-rate versus RIS size $N$ at $P=15$~dBm. NOMA dominates
  OMA at every aperture; both scale with $N$.}
  \label{fig:rateN}
\end{figure}

Finally, Fig.~\ref{fig:rateP} sweeps the transmit power. Panel~(a) establishes
the ordering NOMA-RIS $>$ NOMA $>$ OMA-RIS $>$ OMA across the whole range:
NOMA's full-time superposition outperforms OMA's orthogonal slots, and the RIS
adds a consistent margin on top of each access scheme (e.g.\ a $\sim\!27\%$ RIS
gain for NOMA at $P=15$~dBm). Panel~(b) shows the corresponding QoS-violation
rate: NOMA-RIS reaches near-zero violation already at low power, whereas the
no-RIS and OMA schemes need substantially more power to become feasible,
quantifying the role of the RIS in jointly meeting the sensing and communication
QoS.

\begin{figure}[t]
  \centering
  \includegraphics[width=\columnwidth]{fig_rate_P.eps}
  \caption{(a)~Weighted sum-rate and (b)~joint QoS-violation rate versus
  transmit power for the four schemes.}
  \label{fig:rateP}
\end{figure}

Finally, two practical points follow. On \emph{runtime}, the online stage is a
single forward pass of $f_{\bm\theta}$ plus $K{=}5$ light inner steps---orders
of magnitude below per-instance AO---so the design fits the short coherence
times that motivate learning-based ISAC. On \emph{constraint granularity},
replacing the sum-rate floor \eqref{eq:qosc} with per-user penalties
($\lambda_n,\lambda_f{>}0$ in~\eqref{eq:sampleloss}) lets the operator dial the
near/far fairness trade-off, shifting the learned split toward the weaker user
at a modest cost in aggregate rate and with no change to the pipeline.

%==============================================================================
\section{Conclusion}
A MISO RIS-assisted NOMA-ISAC downlink was investigated with joint power and
phase design. The revised formulation combines an aggregate NOMA-cluster QoS
constraint with individual near/far-user floors and a separate sensing floor,
which preserves cluster efficiency while preventing a strong user from masking
a weak one. The implemented solver uses a compact softmax power network and
per-channel RIS-phase adaptation, followed by a first-order meta-update. The
paper now reports the actual channel distances, fading/path-loss parameters,
power/RIS settings, convergence limitations, and online prediction complexity.
The original result figures and their surrounding discussion are retained and the figure files are provided in EPS form for submission.

%==============================================================================
\begin{thebibliography}{1}
\bibitem{liu2022isac}
F.~Liu \emph{et al.}, ``Integrated sensing and communications: Toward
dual-functional wireless networks for 6G and beyond,''
\emph{IEEE J. Sel. Areas Commun.}, vol.~40, no.~6, pp.~1728--1767, Jun.~2022.

\bibitem{deng2025mmwave}
K.~Deng, X.~Wang, H.~Xu, Q.~Song, R.~Zhang, and Y.~Qian, ``Balancing
communication and sensing hybrid beamforming with rank-1 optimization for
mmWave ISAC systems,'' \emph{IEEE Wireless Commun. Lett.}, vol.~14, no.~9,
pp.~2738--2742, Sep.~2025.

\bibitem{wu2020ris}
Q.~Wu and R.~Zhang, ``Towards smart and reconfigurable environment: Intelligent
reflecting surface aided wireless network,'' \emph{IEEE Commun. Mag.}, vol.~58,
no.~1, pp.~106--112, Jan.~2020.

\bibitem{wang2025hybrid}
L.~Wang, L.~F. Abanto-Leon, and A.~Asadi, ``Joint hybrid beamforming and RIS
phase shift design for RIS-enabled mmWave ISAC system,'' \emph{IEEE Trans. Veh.
Technol.}, vol.~74, no.~6, pp.~9149--9164, Jun.~2025.

\bibitem{lyu2024crb}
W.~Lyu \emph{et al.}, ``CRB minimization for RIS-aided mmWave integrated sensing
and communications,'' \emph{IEEE Internet Things J.}, vol.~11, no.~10,
pp.~18381--18393, May~2024.

\bibitem{ke2025rsma}
H.~Ke, J.~Xu, W.~Xu, C.~Ding, and D.~W.~K. Ng, ``Rate splitting and beamforming
design for RIS-assisted RSMA-enhanced ISAC networks,'' \emph{IEEE Wireless
Commun. Lett.}, vol.~14, no.~6, pp.~1738--1742, Jun.~2025.

\bibitem{zou2023ml}
Y.~Zou, Y.~Liu, X.~Mu, X.~Zhang, Y.~Liu, and C.~Yuen, ``Machine learning in
RIS-assisted NOMA IoT networks,'' \emph{IEEE Internet Things J.}, vol.~10,
no.~22, pp.~19427--19440, Nov.~2023.

\bibitem{wei2026ris}
T.~Wei \emph{et al.}, ``Machine learning-driven joint RIS partitioning,
beamforming, and phase shifts design for mmWave ISAC systems,'' \emph{IEEE
Wireless Commun. Lett.}, vol.~15, 2026.

\bibitem{maml}
C.~Finn, P.~Abbeel, and S.~Levine, ``Model-agnostic meta-learning for fast
adaptation of deep networks,'' in \emph{Proc. ICML}, 2017, pp.~1126--1135.

\bibitem{fallah2020maml}
A.~Fallah, A.~Mokhtari, and A.~Ozdaglar, ``On the convergence theory of
gradient-based model-agnostic meta-learning algorithms,'' in \emph{Proc.
AISTATS}, 2020, pp.~1082--1092.

\bibitem{ji2022maml}
K.~Ji, J.~Yang, and Y.~Liang, ``Theoretical convergence of multi-step
model-agnostic meta-learning,'' \emph{J. Mach. Learn. Res.}, vol.~23, pp.~1--41,
2022.

\bibitem{sun2024noma}
Y.~Sun \emph{et al.}, ``On the study of NOMA-assisted integrated sensing and
communication networks,'' \emph{IEEE Trans. Commun.}, vol.~72, no.~11,
pp.~7278--7293, Nov.~2024.
\end{thebibliography}

\end{document}
